\documentclass[conference,a4paper]{IEEEtran}

% Escritura mejorada de fórmulas matemáticas
\usepackage{amsmath}

% Inserción de gráficos
\usepackage{graphicx}

% Escritura de pseudocódigo
\usepackage[kw]{pseudo}

% Escritura mejorada de tablas
\usepackage{booktabs}

% Escritura mejorada de citas bibliográficas
\usepackage{cite}


% Macros traducidas
\def\contentsname{Índice general}
\def\listfigurename{Índice de figuras}
\def\listtablename{Índice de tablas}
\def\refname{Referencias}
\def\indexname{Índice alfabético}
\def\figurename{Fig.}
\def\tablename{TABLA}
\def\partname{Parte}
\def\appendixname{Apéndice}
\def\abstractname{Resumen}
% IEEE specific names
\def\IEEEkeywordsname{Palabras clave}
\def\IEEEproofname{Demostración}


\begin{document}

\title{Título del trabajo (elegir uno original)}

\author{
  \IEEEauthorblockN{Sandra Silva Amaya}
  \IEEEauthorblockA{
    \textit{Dpto. Ciencias de la Computación e Inteligencia Artificial}\\
    \textit{Universidad de Sevilla}\\
    Sevilla, España\\
    sansilama@alum.us.es}
  
  \and
  
  \IEEEauthorblockN{Nuria Martín Sánchez}
  \IEEEauthorblockA{
    \textit{Dpto. Ciencias de la Computación e Inteligencia Artificial}\\
    \textit{Universidad de Sevilla}\\
    Sevilla, España\\
    Correos electrónicos UVUS y de contacto (si distinto)}
}

\maketitle


% Resumen
\begin{abstract}
 El objetivo de este trabajo es desarrollar y comparar dos sistemas de 
 recuperación de información. El primero de ellos se basa en el modelo de 
 recuperación booleano, mientras que el segundo admite consultas de texto libre, 
 empleando el modelo de ponderación TF-IDF y la similitud del coseno para 
 devolver una lista ordenada de resultados. Para implementar ambos sistemas, 
 se llevó a cabo un preprocesamiento del lenguaje natural sobre el corpus 
 de prueba mediante tokenización, eliminación de palabras vacías (stop words) y 
 extracción de raíces (stemming), apoyándose en el uso de librerías como NLTK y Whoosh.

 La evaluación de ambos sistemas se realizó utilizando métricas de precisión y 
 exhaustividad (recall). El sistema booleano alcanzó una precisión de 0.998 y 
 un recall de 1.0, mientras que el modelo TF-IDF obtuvo un Mean Average 
 Precision (MAP) de 0.7248. No obstante, se determinó que estos resultados no 
 implican necesariamente una superioridad del modelo booleano en escenarios reales. 
 Debido a que el corpus de prueba y las necesidades de información fueron diseñados 
 por las mismas autoras del proyecto, las consultas booleanas se beneficiaron de un 
 conocimiento previo del vocabulario indexado. Teniendo en cuenta este sesgo 
 conceptual, se concluye que el modelo TF-IDF presentaría un rendimiento más 
 robusto y realista frente a una colección de prueba evaluada por usuarios 
 externos que no conozcan de antemano el contenido exacto de los documentos.

\end{abstract}


% Palabras claves
\begin{IEEEkeywords}
  Recuperación de información, modelo booleano, TF-IDF, 
  procesamiento del lenguaje natural, evaluación de sistemas de búsqueda
\end{IEEEkeywords}


\section{Introducción}

Como campo académico de estudio, la recuperación de información (RI) se define 
como el proceso de buscar y encontrar material de naturaleza no estructurada, 
predominantemente texto, que satisfaga una necesidad de información dentro de 
grandes colecciones de documentos~\cite{b1}. En la actualidad, esta disciplina ha 
superado con creces los entornos tradicionales de bases de datos relacionales, 
convirtiéndose en la tecnología subyacente fundamental en motores de búsqueda web, 
sistemas de correo electrónico y plataformas de recomendación de contenido 
multimedia~\cite{b2}. Para gestionar este volumen de información textual, se han 
desarrollado históricamente diversos modelos matemáticos de recuperación, los 
cuales varían en complejidad y enfoque lógico, abarcando desde las coincidencias 
lógicas exactas hasta los sistemas basados en espacios vectoriales.

En el ámbito del entretenimiento y el cine, la correcta interpretación de las consultas 
de los usuarios resulta crítica para garantizar una buena experiencia de uso. Sin 
embargo, las colecciones de documentos que describen este contenido suelen adolecer 
de ruido conceptual y variaciones lingüísticas. Para mitigar esta problemática, la 
aplicación de técnicas de preprocesamiento de lenguaje natural (PLN), tales como 
la tokenización, la eliminación de palabras vacías (stop words) y la reducción de 
términos a sus raíces morfológicas (stemming), resulta indispensable para 
homogeneizar los textos antes de su indexación~\cite{b3}. Una vez estructurado el 
corpus, surge el desafío de evaluar la eficacia de los motores de búsqueda mediante 
métricas estandarizadas como la precisión, la exhaustividad (recall) y el 
promedio de precisión media (Mean Average Precision o MAP).

El objetivo principal de este trabajo es el diseño, la implementación y la evaluación 
comparativa de dos aproximaciones clásicas en la recuperación de información. El 
primer sistema propuesto se fundamenta en el modelo de recuperación Booleano, 
implementado mediante la biblioteca Whoosh, el cual opera bajo una lógica de 
correspondencia estricta. El segundo sistema consiste en un motor de consultas en 
texto libre basado en el modelo de espacio vectorial, utilizando la ponderación TF-IDF 
y el cálculo de la similitud del coseno para establecer un ranking ordenado de 
relevancia. Ambos enfoques son validados empleando un corpus propio de sinopsis y 
reseñas cinematográficas, confeccionado junto con un banco de veinte necesidades de 
información y sus respectivos juicios de relevancia. A través de este escenario, se 
persigue analizar las limitaciones prácticas de ambos modelos, prestando especial 
atención a fenómenos como el sobreajuste y el sesgo de autoría en la evaluación.

El resto de este documento se estructura de la siguiente manera: la Sección II 
detalla las características del corpus cinematográfico y los procesos de 
preprocesamiento de lenguaje natural aplicados; la Sección III describe la 
arquitectura matemática y el desarrollo de los modelos de recuperación Booleano y 
Vectorial; la Sección IV expone la metodología de evaluación y analiza 
comparativamente las métricas obtenidas por ambos sistemas. Finalmente, la Sección V 
recoge las conclusiones principales derivadas de este estudio y propone posibles líneas 
de trabajo futuro.


\section{Preliminares}

En esta sección se hace una breve introducción de las técnicas empleadas y
también trabajos relacionados, si los hay.


\subsection{Métodos empleados}

Describir aquí los métodos y técnicas empleadas (búsqueda en espacio de
estados, algoritmos genéticos, redes bayesianas, técnicas de clasificación,
redes neuronales, etc.). Si es necesario, separarlos en distintas subsecciones
dentro de Antecedentes.

Se pueden usar listas por puntos como sigue:
\begin{itemize}
\item Un punto: esto es un ejemplo de una lista.
\item Otro punto.
\end{itemize}

Por último, se debe hacer un uso correcto de las referencias bibliográficas,
para que el lector pueda acceder a más información~\cite{b2}. Todas las
referencias al final del documento deben ser citadas al menos una vez.


\subsection{Trabajo Relacionado}

Se puede realizar un recorrido en la literatura sobre trabajos anteriores que
estén relacionados y que sea por tanto interesante comentar aquí. Por supuesto,
añadir las referencias bibliográficas correspondientes.


\section{Metodología}

Esta sección se dedica a la descripción del método implementado en el trabajo.
Esta parte es la correspondiente a lo realmente desarrollado en el trabajo, y
se puede emplear pseudocódigo (nunca código), esquemas, tablas, etc.

\begin{figure}
  \centering
  \includegraphics{ejemplo}
  \caption{Ejemplo de un pie de figura. Imagen con derechos Creative Commons}
  \label{fig:ejemplo}
\end{figure}

A continuación, un ejemplo de uso de listas numeradas:
\begin{enumerate}
\item\label{item:dos-alumnos} \textit{Trabajos con dos alumnos:} poner nombre y
  apellidos completos de cada uno, y correos electrónicos de contacto (a ser
  posible de la Universidad de Sevilla). El orden de los alumnos se fijará por
  orden alfabético según los apellidos.
\item \textit{Trabajo con un autor:} cambiar la cabecera de la siguiente manera
  \begin{enumerate}
  \item \textit{Una sola columna:} solo se debe especificar un alumno.
  \item \textit{Información a añadir:} la misma que la especificada en el
    punto~\ref{item:dos-alumnos}.
  \end{enumerate}
\end{enumerate}

Las figuras se deben mencionar en el texto, como la
\figurename~\ref{fig:ejemplo}. También se pueden añadir ecuaciones, como la
ecuación~\eqref{eq:ejemplo}.

\begin{equation}
  \label{eq:ejemplo}
  a + b = \gamma
\end{equation}

Un ejemplo de pseudocódigo se puede observar en la
\figurename~\ref{pcd:mergesort}.

\begin{figure}
  \begin{pseudo}*
    \hd{\fn{mergesort}}(V) \\*
    \multicolumn{2}{l}{\textbf{Entrada}: un vector \( V \)} \\*
    \multicolumn{2}{l}{\textbf{Salida}: un vector con los elementos de \( V \)
      en orden} \\
    si \( V \) \textnormal{es unitario} entonces \\+
    devolver \( V \) \\-
    si no entonces \\+
    \( V_{1} \leftarrow \textnormal{primera mitad de } V\) \\
    \( V_{2} \leftarrow \textnormal{segunda mitad de } V\) \\
    \( V_{1} \leftarrow \pr{mergesort}(V_{1}) \) \\
    \( V_{2} \leftarrow \pr{mergesort}(V_{2}) \) \\
    devolver \fn{mezcla}(V_{1}, V_{2})
  \end{pseudo}
  
  \begin{pseudo}*
    \hd{\fn{mezcla}}(V_{1}, V_{2}) \\*
    \multicolumn{2}{l}{%
      \textbf{Entrada}: dos vectores \( V_{1} \) y \( V_{2} \) ordenados
    } \\*
    \multicolumn{2}{l}{%
      \textbf{Salida}: un vector con los elementos de \( V_{1} \) y \( V_{2} \)
      en orden
    } \\
    si \( V_{1} \) \textnormal{no tiene elementos} entonces \\+
    devolver \( V_{2} \) \\-
    si no si \( V_{2} \) \textnormal{no tiene elementos} entonces \\+
    devolver \( V_{1} \) \\-
    si no entonces \\+
    \( x_{1} \leftarrow \textnormal{primer elemento de } V_{1} \) \\
    \( x_{2} \leftarrow \textnormal{primer elemento de } V_{2} \) \\
    si \( x_{1} \leq x_{2} \) entonces \\+
    \( x \leftarrow x_{1} \) \\
    \textnormal{quitar el primer elemento de} \( V_{1} \) \\-
    si no entonces \\+
    \( x \leftarrow x_{2} \) \\
    \textnormal{quitar el primer elemento de} \( V_{2} \) \\-
    \( V \leftarrow \fn{mezcla}(V_{1}, V_{2}) \) \\
    \textnormal{añadir} \( x \)
    \textnormal{como primer elemento de} \( V \) \\
    devolver \( V \)
  \end{pseudo}
  \caption{Algoritmo de ordenación \texttt{MergeSort}}
  \label{pcd:mergesort}
\end{figure}


\section{Resultados}

En esta sección se detallarán tanto los experimentos realizados como los
resultados conseguidos:
\begin{itemize}
\item Los experimentos realizados, indicando razonadamente la configuración
  empleada, qué se quiere determinar, y como se ha medido.
\item Los resultados obtenidos en cada experimento, explicando en cada caso lo
  que se ha conseguido.
\item Análisis de los resultados, haciendo comparativas y obteniendo
  conclusiones.
\end{itemize}

Se pueden hacer uso de tablas, como el ejemplo de la tabla~\ref{tab:ejemplo}.

\begin{table}
  \caption{Ejemplo de tabla}
  \label{tab:ejemplo}
  \centering
  \begin{tabular}{ccc}
    \toprule
    A & B & C \\
    \midrule
    1 & 2 & 3 \\
    4 & 5 & 6 \\
    \bottomrule
  \end{tabular}
\end{table}


\section{Conclusiones}

Finalmente, se dedica la última sección para indicar las conclusiones obtenidas
del trabajo. Se puede dedicar un párrafo para realizar un resumen sucinto del
trabajo, con los experimentos y resultados. Seguidamente, uno o dos párrafos
con conclusiones. Se suele dedicar un párrafo final con ideas de mejora y
trabajo futuro.


\begin{thebibliography}{00}
\bibitem{b1} C. D. Manning, P. Raghavan, y H. Schütze, \emph{Introduction to Information Retrieval}. Cambridge, Reino Unido: Cambridge University Press, 2008.
\bibitem{b2} Á. Romero Jiménez, ``Propuesta de trabajo práctico de la asignatura Inteligencia Artificial: Recuperación de la información,'' Dpto. de Ciencias de la Computación e Inteligencia Artificial, Universidad de Sevilla, Curso 2025/2026.
\bibitem{b3} S. Bird, E. Klein, y E. Loper, \emph{Natural Language Processing with Python}. Sebastopol, CA, EE. UU.: O'Reilly Media, 2009.
\end{thebibliography}


\end{document}
